\documentclass[12pt, a4paper]{article}
% --- 1. Cấu hình Tiếng Việt và Font chữ chuẩn ---
\usepackage[utf8]{inputenc}
\usepackage[T5]{fontenc}
\usepackage[vietnamese]{babel}
\usepackage{mathptmx} % Font Times New Roman đồng bộ cả chữ và công thức toán, không gây xung đột gói

\makeatletter
\renewcommand{\normalsize}{\@setfontsize\normalsize{13pt}{16pt}}
\makeatother
\normalsize

% --- 2. Gói Toán học & Ký hiệu bổ trợ ---
\usepackage{amsmath, amssymb, amsfonts, mathrsfs, amsthm}
\usepackage{graphicx}
\usepackage{fontawesome}
\usepackage{booktabs, tabularx, array, multirow}
\usepackage{enumitem}

% --- 3. Đồ họa TikZ, Bảng biến thiên & Hình không gian ---
\usepackage{tikz}
\usetikzlibrary{calc, intersections, angles, quotes, patterns, arrows.meta, 3d}
\usepackage{pgfplots}
\pgfplotsset{compat=1.18}
\usepackage{tkz-euclide}
\usepackage{tkz-tab}

\renewcommand{\vec}[1]{\overrightarrow{#1}}

% --- 4. Hộp sư phạm (Tcolorbox) ---
\usepackage[many]{tcolorbox}
\tcbuselibrary{skins, breakable}

% --- 5. Căn lề chuẩn văn bản giáo án ---
\usepackage{geometry}
\geometry{
    a4paper,
    left=25mm,
    right=15mm,
    top=20mm,
    bottom=20mm
}

% --- 6. Header / Footer chuẩn hóa ---
\usepackage{fancyhdr}
\usepackage{lastpage}
\pagestyle{fancy}
\fancyhf{}

\fancyhead[L]{\small \textit{Kế hoạch bài dạy Toán 11}}
\fancyhead[R]{\textbf{Trang \thepage/\pageref{LastPage}}}
\fancyfoot[L]{\small \textit{Tổ Toán - Tin}}
\fancyfoot[R]{\small \textit{Trường THCS\&THPT Hồng Vân}}
\renewcommand{\headrulewidth}{0.4pt}
\renewcommand{\footrulewidth}{0.4pt}

% --- 7. Giãn dòng & Giãn đoạn theo chuẩn Nghị định 30/2020/NĐ-CP ---
\usepackage{setspace}
\setstretch{1.15}
\setlength{\parindent}{1.0cm}
\setlength{\parskip}{5pt}

% Định nghĩa hộp sư phạm chuyên dụng
\newtcolorbox{pedabox}[2][]{
    enhanced,
    breakable,
    colback=blue!3!white,
    colframe=blue!65!black,
    fonttitle=\bfseries\small,
    title={#2},
    arc=2mm,
    boxrule=0.6pt,
    left=3mm, right=3mm, top=2mm, bottom=2mm,
    #1
}

\newtcolorbox{aibox}[2][]{
    enhanced,
    breakable,
    colback=teal!4!white,
    colframe=teal!70!black,
    fonttitle=\bfseries\small,
    title={#2},
    arc=2mm,
    boxrule=0.6pt,
    left=3mm, right=3mm, top=2mm, bottom=2mm,
    #1
}

\newtcolorbox{scaffoldbox}[2][]{
    enhanced,
    breakable,
    colback=orange!4!white,
    colframe=orange!80!black,
    fonttitle=\bfseries\small,
    title={#2},
    arc=2mm,
    boxrule=0.6pt,
    left=3mm, right=3mm, top=2mm, bottom=2mm,
    #1
}

\begin{document}

\begin{center}
    {\fontsize{14pt}{17pt}\selectfont \textbf{KẾ HOẠCH BÀI DẠY MÔN TOÁN LỚP 11}}\\[6pt]
    {\fontsize{16pt}{19pt}\selectfont \textbf{BÀI 16. GIỚI HẠN CỦA HÀM SỐ}}\\[4pt]
    {\fontsize{12pt}{15pt}\selectfont \textbf{Môn học: Toán 11} \ \textbar \ \textbf{Thời lượng: 2 tiết (Tiết 42, 43 theo PPCT | Tuần 14, 15)}}
\end{center}

\vspace{5pt}

\section*{I. MỤC TIÊU}

\subsection*{1. Kiến thức, Năng lực}

\begin{itemize}[leftmargin=1.2cm]
    \item \textbf{Yêu cầu cần đạt (Nguyên văn CT GDPT 2018 Toán 11):}
    \begin{itemize}[leftmargin=0.6cm]
        \item Nhận biết được khái niệm giới hạn hữu hạn của hàm số, giới hạn hữu hạn một phía của hàm số tại một điểm.
        \item Nhận biết được khái niệm giới hạn hữu hạn của hàm số tại vô cực và mô tả được một số giới hạn cơ bản như: $\lim_{x \to +\infty} \dfrac{c}{x^k} = 0$, $\lim_{x \to -\infty} \dfrac{c}{x^k} = 0$ ($k \in \mathbb{N}^*$, $c$ là hằng số).
        \item Nhận biết được khái niệm giới hạn vô cực (một phía) của hàm số tại một điểm.
        \item Tính được một số giới hạn hàm số bằng cách vận dụng các phép toán trên giới hạn hàm số (tổng, hiệu, tích, thương, căn thức) và khử các dạng vô định cơ bản (dạng $\dfrac{0}{0}$).
        \item Giải quyết được một số vấn đề thực tiễn gắn với giới hạn hàm số (chi phí sản xuất trung bình tiệm cận, vận tốc tức thời, tải trọng tới hạn của kết cấu).
    \end{itemize}

    \item \textbf{Năng lực Toán học đặc thù:}
    \begin{itemize}[leftmargin=0.6cm]
        \item \textit{Năng lực tư duy và lập luận toán học:} Phân biệt rành mạch giữa quá trình biến thiên của đối số $x \to x_0$ (tiệm cận) và giá trị của hàm số tại điểm đó $f(x_0)$ (có thể không xác định); hiểu sâu sắc mối quan hệ hai chiều giữa giới hạn hai phía và giới hạn một phía: $\lim_{x \to x_0} f(x) = L \iff \lim_{x \to x_0^+} f(x) = \lim_{x \to x_0^-} f(x) = L$.
        \item \textit{Năng lực mô hình hóa toán học:} Vận dụng giới hạn tại vô cực để mô tả xu hướng dài hạn của các quá trình kinh tế - kỹ thuật (hàm chi phí trung bình tiệm cận khi quy mô sản xuất tăng vô hạn) và giới hạn vô cực để mô phỏng tải trọng tới hạn trong cơ học.
        \item \textit{Năng lực giải quyết vấn đề toán học:} Nắm vững thuật toán xử lý các dạng bài: phương pháp thay số trực tiếp, phương pháp phân tích nhân tử để khử dạng vô định $\dfrac{0}{0}$ ở biểu thức hữu tỉ, phương pháp nhân lượng liên hợp khử dạng vô định chứa căn thức, và kỹ thuật chia lũy thừa bậc cao nhất khi $x \to \pm\infty$.
        \item \textit{Năng lực giao tiếp toán học:} Trình bày lời giải giải tích chuẩn xác; sử dụng chính xác hệ thống ký hiệu $\lim_{x \to x_0^+}, \lim_{x \to x_0^-}, \lim_{x \to +\infty}, \lim_{x \to -\infty}$; vẽ và đọc hiểu hình ảnh đồ thị hàm số có điểm gián đoạn hoặc đường tiệm cận.
        \item \textit{Năng lực sử dụng công cụ, phương tiện học toán:} Khai thác phần mềm hình học trực quan Desmos / GeoGebra để quan sát đồ thị tiếp cận điểm giới hạn; thao tác thành thạo MTCT Casio fx-580VN X để tính giới hạn một phía và hai phía.
    \end{itemize}

    \item \textbf{Năng lực số (Theo Thông tư 02/2025/TT-BGDĐT):}
    \begin{itemize}[leftmargin=0.6cm]
        \item \textit{Mã 1.1NC1b (Khai thác học liệu số và công cụ mô phỏng):} Học sinh biết cách truy cập và sử dụng phần mềm đồ họa hàm số Desmos / GeoGebra để quan sát hành vi của đồ thị hàm số khi đối số $x$ tiến rất gần đến điểm kỳ dị $x_0$.
        \item \textit{Mã 5.2NC1b (Giải quyết vấn đề toán học với công cụ tính toán số):} Học sinh áp dụng thành thạo máy tính cầm tay (MTCT Casio fx-580VN X) sử dụng chức năng phím CALC với giá trị $x = x_0 + 10^{-9}$ (đối với giới hạn phải $x \to x_0^+$) và $x = x_0 - 10^{-9}$ (đối với giới hạn trái $x \to x_0^-$) để dự đoán chính xác giới hạn một phía và kiểm tra dạng vô định.
    \end{itemize}

    \item \textbf{Năng lực AI (Theo Quyết định 2422/QĐ-BGDĐT):}
    \begin{itemize}[leftmargin=0.6cm]
        \item \textit{Mã 11.C3.1 (Kỹ thuật đặt câu lệnh Prompt tương tác với AI):} Học sinh biết cách xây dựng prompt có cấu trúc rõ ràng yêu cầu trợ lý AI giải thích trực quan sự khác biệt giữa khái niệm $x \to x_0$ và $x = x_0$, đồng thời nhận diện được trường hợp AI đưa ra lời giải sai khi hàm số không tồn tại giới hạn.
        \item \textit{Mã 11.C2.1 (Nhận thức ứng dụng chuyên sâu của AI trong thực tiễn):} Học sinh nhận biết cách các hệ thống AI mô phỏng vật lý và tính toán kỹ thuật số sử dụng giới hạn vô cực để xác định ngưỡng tải trọng tới hạn, ngăn ngừa nguy cơ sụp đổ công trình xây dựng.
    \end{itemize}

    \item \textbf{Năng lực chung:}
    \begin{itemize}[leftmargin=0.6cm]
        \item \textit{Tự chủ và tự học:} Tự giác nghiên cứu các ví dụ mẫu trong SGK; chủ động rà soát các bước nhân liên hợp và phân tích đa thức thành nhân tử.
        \item \textit{Giao tiếp và hợp tác:} Tích cực phối hợp với bạn học trong hoạt động nhóm, thảo luận sôi nổi về sự tồn tại của giới hạn một phía.
    \end{itemize}
\end{itemize}

\subsection*{2. Phẩm chất}

\begin{itemize}[leftmargin=1.2cm]
    \item \textbf{Chăm chỉ:} Kiên nhẫn thực hiện đầy đủ các bước phân tích nhân tử và rút gọn biểu thức đại số phức tạp.
    \item \textbf{Trung thực:} Tự giác tính toán, nghiêm túc kiểm tra kết quả bằng MTCT và chỉ ra các sơ hở trong câu trả lời của AI.
    \item \textbf{Trách nhiệm:} Hợp tác tích cực trong nhóm, chủ động hỗ trợ học sinh trung bình - yếu làm chủ kỹ năng xét dấu phân thức.
\end{itemize}

\section*{II. THIẾT BỊ DẠY HỌC VÀ HỌC LIỆU}

\begin{itemize}[leftmargin=1.2cm]
    \item \textbf{Giáo viên:}
    \begin{itemize}[leftmargin=0.6cm]
        \item Kế hoạch bài dạy, bài giảng điện tử PowerPoint/Canva tương tác minh họa sự tiếp cận điểm giới hạn.
        \item Mô phỏng động trên Desmos: Hàm số $f(x) = \dfrac{x^2 - 1}{x - 1}$ và $g(x) = \dfrac{1}{x - 2}$ với điểm chạy $M(x; f(x))$ tiến sát đến điểm giới hạn.
        \item Phiếu học tập phân tầng bậc thang (Worksheet), bảng Rubric đánh giá năng lực học sinh.
        \item Máy tính cầm tay Casio fx-580VN X kết nối phần mềm giả lập chiếu lên màn hình lớp học.
    \end{itemize}
    \item \textbf{Học sinh:}
    \begin{itemize}[leftmargin=0.6cm]
        \item Sách giáo khoa Toán 11, vở ghi bài, thước kẻ, bút viết.
        \item Máy tính cầm tay Casio fx-580VN X đã sẵn sàng cho thao tác lệnh CALC.
    \end{itemize}
\end{itemize}

\section*{III. TIẾN TRÌNH DẠY HỌC (BAO QUÁT TRỌN VẸN 2 TIẾT)}

\subsection*{TIẾT 1: KHÁI NIỆM GIỚI HẠN HÀM SỐ, GIỚI HẠN MỘT PHÍA VÀ GIỚI HẠN TẠI VÔ CỰC (TIẾT 42 THEO PPCT)}

\subsubsection*{1. Hoạt động 1: Mở đầu (Khởi động) (5 phút)}

\noindent \textbf{a) Mục tiêu:}
Tạo tình huống mâu thuẫn nhận thức khi tính giá trị hàm số tại điểm làm mẫu thức bằng 0, từ đó dẫn dắt tự nhiên vào khái niệm giới hạn hàm số tại một điểm.

\noindent \textbf{b) Nội dung: Tình huống hàm số có điểm thủng tại tập xác định}
\begin{pedabox}{Khởi động: Giá trị tại điểm và xu hướng khi tiến sát tới điểm}
Xét hàm số $f(x) = \dfrac{x^2 - 1}{x - 1}$ xác định trên tập hợp $\mathbb{R} \setminus \{1\}$.
\begin{enumerate}[leftmargin=0.6cm]
    \item Ta có tính được giá trị $f(1)$ hay không? Vì sao?
    \item Lập bảng tính giá trị $f(x)$ khi cho $x$ nhận các giá trị rất gần số $1$:
    \begin{center}
    \small
    \begin{tabular}{|c|c|c|c|c|c|c|}
    \hline
    $x$ & $0{,}9$ & $0{,}99$ & $0{,}999$ & $1{,}001$ & $1{,}01$ & $1{,}1$ \\ \hline
    $f(x)$ & $1{,}9$ & $1{,}99$ & $1{,}999$ & $2{,}001$ & $2{,}01$ & $2{,}1$ \\ \hline
    \end{tabular}
    \end{center}
    \item Khi $x$ tiến ngày càng gần đến số $1$ (nhưng $x \neq 1$), giá trị $f(x)$ tiến dần về số thực nào?
\end{enumerate}
\end{pedabox}

\noindent \textbf{c) Sản phẩm: Câu trả lời và nhận xét của học sinh}
\begin{itemize}[leftmargin=0.6cm]
    \item Hàm số không xác định tại $x = 1$ nên không tồn tại giá trị $f(1)$ (phép chia cho 0 không có nghĩa).
    \item Với mọi $x \neq 1$, ta có thể rút gọn: $f(x) = \dfrac{(x-1)(x+1)}{x-1} = x + 1$.
    \item Khi $x$ tiến rất sát đến $1$ (từ bên trái hay bên phải), giá trị của $f(x) = x + 1$ tiến sát về số $2$.
    \item Học sinh kết luận: Tuy không có $f(1)$ nhưng khi $x$ tiến dần về $1$ thì $f(x)$ tiến dần về $2$. Ta nói giới hạn của $f(x)$ khi $x$ dần tới $1$ bằng $2$.
\end{itemize}

\noindent \textbf{d) Tổ chức thực hiện:}
\begin{itemize}[leftmargin=0.8cm]
    \item \textit{Chuyển giao nhiệm vụ:} Giáo viên trình chiếu bảng số liệu; yêu cầu học sinh thảo luận cặp đôi trả lời 3 câu hỏi trong 2 phút.
    \item \textit{Thực hiện nhiệm vụ:} Học sinh quan sát bảng số liệu, bấm máy tính kiểm tra một vài điểm; nhận thấy quy luật $f(x) \to 2$.
    \item \textit{Báo cáo, thảo luận:} Giáo viên gọi đại diện 1 bàn trả lời; cả lớp thống nhất nhận định.
    \item \textit{Kết luận, nhận định:} Giáo viên nhấn mạnh: Giới hạn hàm số tại điểm $x_0$ phản ánh \textbf{xu thế biến thiên} của hàm số khi đối số $x$ đến rất gần $x_0$, hoàn toàn độc lập với việc hàm số có xác định tại $x_0$ hay không.
\end{itemize}

\subsubsection*{2. Hoạt động 2: Hình thành kiến thức (20 phút)}

\noindent \textbf{a) Mục tiêu:}
\begin{itemize}[leftmargin=0.6cm]
    \item Nắm vững định nghĩa giới hạn hữu hạn của hàm số tại một điểm, định lý về các phép toán trên giới hạn.
    \item Hiểu rõ khái niệm giới hạn một phía và định lý điều kiện tồn tại giới hạn hai phía.
    \item Nhận biết giới hạn hữu hạn tại vô cực và các giới hạn cơ bản $\lim_{x \to \pm\infty} \dfrac{c}{x^k} = 0$.
    \item Tích hợp Năng lực số 1.1NC1b (Desmos) và Năng lực AI 11.C3.1 (kỹ thuật đặt câu lệnh prompt).
\end{itemize}

\noindent \textbf{b) Nội dung: Định nghĩa, phép toán, giới hạn một phía và giới hạn tại vô cực}

\begin{pedabox}{1. Khái niệm giới hạn hữu hạn của hàm số tại một điểm}
Cho khoảng $K$ chứa điểm $x_0$ và hàm số $y = f(x)$ xác định trên $K$ hoặc trên $K \setminus \{x_0\}$.
\begin{itemize}[leftmargin=0.6cm]
    \item Ta nói hàm số $y = f(x)$ có \textbf{giới hạn hữu hạn là số $L$} khi $x$ dần tới $x_0$ nếu với dãy số $(x_n)$ bất kì, $x_n \in K \setminus \{x_0\}$ và $x_n \to x_0$, ta đều có $f(x_n) \to L$.
    \item Kí hiệu: $\lim_{x \to x_0} f(x) = L$ hay $f(x) \to L$ khi $x \to x_0$.
\end{itemize}
\end{pedabox}

\noindent \textbf{Hình vẽ minh họa đồ thị có điểm thủng tại $(1; 2)$:}
\begin{center}
\begin{tikzpicture}[scale=1.1]
    \draw[->, thick] (-1,0) -- (3.5,0) node[right] {$x$};
    \draw[->, thick] (0,-0.5) -- (0,3.5) node[above] {$y$};
    \draw[dashed] (1,0) node[below] {$1$} -- (1,2) -- (0,2) node[left] {$2$};
    % Vẽ đường thẳng y = x + 1 bỏ điểm (1,2)
    \draw[thick, blue] (-0.5,0.5) -- (0.95,1.95);
    \draw[thick, blue] (1.05,2.05) -- (2.3,3.3) node[right] {$y = \dfrac{x^2 - 1}{x - 1}$};
    \draw[thick, red] (1,2) circle (2.5pt); % Điểm thủng
    \node[above left, red] at (1,2) {\small Điểm thủng $(1; 2)$};
    % Mũi tên tiến tới 1
    \draw[->, ultra thick, teal] (0.2,-0.25) -- (0.9,-0.25) node[midway, below] {\small $x \to 1^-$};
    \draw[->, ultra thick, teal] (1.8,-0.25) -- (1.1,-0.25) node[midway, below] {\small $x \to 1^+$};
\end{tikzpicture}
\end{center}

\begin{pedabox}{2. Các phép toán trên giới hạn hữu hạn của hàm số}
Nếu $\lim_{x \to x_0} f(x) = L$ và $\lim_{x \to x_0} g(x) = M$ ($L, M \in \mathbb{R}$) thì:
\begin{enumerate}[leftmargin=0.6cm]
    \item $\lim_{x \to x_0} [f(x) \pm g(x)] = L \pm M$.
    \item $\lim_{x \to x_0} [f(x) \cdot g(x)] = L \cdot M$.
    \item $\lim_{x \to x_0} \dfrac{f(x)}{g(x)} = \dfrac{L}{M}$ (với điều kiện $M \neq 0$).
    \item Nếu $f(x) \ge 0$ với mọi $x \in K \setminus \{x_0\}$ thì $L \ge 0$ và $\lim_{x \to x_0} \sqrt{f(x)} = \sqrt{L}$.
    \item Giới hạn cơ bản: $\lim_{x \to x_0} c = c$ ($c$ là hằng số); $\lim_{x \to x_0} x = x_0$.
\end{enumerate}
\end{pedabox}

\begin{pedabox}{3. Giới hạn một phía}
\begin{itemize}[leftmargin=0.6cm]
    \item \textbf{Giới hạn bên phải:} $\lim_{x \to x_0^+} f(x) = L$ nếu với mọi dãy $(x_n)$ sao cho $x_n > x_0$ và $x_n \to x_0$ thì $f(x_n) \to L$.
    \item \textbf{Giới hạn bên trái:} $\lim_{x \to x_0^-} f(x) = L$ nếu với mọi dãy $(x_n)$ sao cho $x_n < x_0$ và $x_n \to x_0$ thì $f(x_n) \to L$.
    \item \textbf{Định lý then chốt:}
    $$\lim_{x \to x_0} f(x) = L \iff \lim_{x \to x_0^+} f(x) = \lim_{x \to x_0^-} f(x) = L.$$
    Nếu $\lim_{x \to x_0^+} f(x) \neq \lim_{x \to x_0^-} f(x)$ thì \textbf{không tồn tại} giới hạn $\lim_{x \to x_0} f(x)$.
\end{itemize}
\end{pedabox}

\begin{pedabox}{4. Giới hạn hữu hạn của hàm số tại vô cực}
\begin{itemize}[leftmargin=0.6cm]
    \item Ta nói hàm số $y = f(x)$ có giới hạn là $L$ khi $x \to +\infty$ nếu với mọi dãy số $(x_n)$ sao cho $x_n \to +\infty$ thì $f(x_n) \to L$. Kí hiệu: $\lim_{x \to +\infty} f(x) = L$. (Tương tự cho $x \to -\infty$).
    \item \textbf{Giới hạn cơ bản:} Với $k$ là số nguyên dương ($k \in \mathbb{N}^*$) và $c$ là hằng số:
    $$\lim_{x \to +\infty} \dfrac{c}{x^k} = 0; \quad \lim_{x \to -\infty} \dfrac{c}{x^k} = 0.$$
\end{itemize}
\end{pedabox}

\begin{scaffoldbox}{Hộp hỗ trợ sư phạm: Mẹo phân biệt $x \to x_0^+$ và $x \to x_0^-$ cho học sinh TB-Yếu}
Trên trục số định hướng từ trái sang phải:
\begin{itemize}[leftmargin=0.6cm]
    \item $x \to x_0^+$: Điểm $x$ tiến về $x_0$ từ \textbf{bên phải} (các giá trị lớn hơn $x_0$). Ví dụ khi $x \to 2^+$, ta tưởng tượng $x = 2{,}001; 2{,}00001$.
    \item $x \to x_0^-$: Điểm $x$ tiến về $x_0$ từ \textbf{bên trái} (các giá trị nhỏ hơn $x_0$). Ví dụ khi $x \to 2^-$, ta tưởng tượng $x = 1{,}999; 1{,}99999$.
\end{itemize}
\end{scaffoldbox}

\begin{aibox}{Năng lực AI 11.C3.1: Kỹ thuật đặt Prompt tương tác giải thích khái niệm Toán học}
\textbf{Câu lệnh Prompt chuẩn mực mà học sinh có thể gửi cho trợ lý AI:}
\begin{quote}
\textit{«Hãy đóng vai một chuyên gia sư phạm Toán học. Em đang học bài Giới hạn của hàm số lớp 11. Xin hãy giải thích bằng ngôn ngữ trực quan, không dùng công thức phức tạp: Tại sao khi nói $x \to 1$ thì giá trị $x$ không bao giờ bằng $1$? Cho một ví dụ thực tế tương tự trong đời sống để em dễ hình dung.»}
\end{quote}
\textbf{Đánh giá câu trả lời của AI:} AI sẽ lấy ví dụ về việc một chiếc máy bay hạ cánh: độ cao của máy bay tiến sát tới mặt đất ($h \to 0$), nhưng trong quá trình tiếp cận nó chưa chạm hẳn xuống đất. Điều này giúp học sinh khắc sâu bản chất của khái niệm tiệm cận.
\end{aibox}

\noindent \textbf{c) Sản phẩm: Lời giải các ví dụ mẫu}
\textbf{Ví dụ 1 (Giới hạn xác định trực tiếp):} Tính $\lim_{x \to 2} (3x^2 - 2x + 1)$.
\begin{itemize}[leftmargin=0.6cm]
    \item \textit{Lời giải:} Áp dụng định lý về phép toán:
    $$\lim_{x \to 2} (3x^2 - 2x + 1) = 3 \cdot (\lim_{x \to 2} x)^2 - 2 \cdot (\lim_{x \to 2} x) + 1 = 3 \cdot (2)^2 - 2 \cdot (2) + 1 = 12 - 4 + 1 = 9.$$
\end{itemize}

\textbf{Ví dụ 2 (Giới hạn một phía của hàm phân nhánh):} Cho hàm số $f(x) = \begin{cases} 2x + 1 & \text{khi } x \ge 1 \\ x^2 + 2 & \text{khi } x < 1 \end{cases}$. Tìm $\lim_{x \to 1^+} f(x)$, $\lim_{x \to 1^-} f(x)$ và cho biết có tồn tại $\lim_{x \to 1} f(x)$ hay không?
\begin{itemize}[leftmargin=0.6cm]
    \item \textit{Lời giải:}
    \begin{itemize}
        \item Khi $x \to 1^+$ thì $x > 1$, do đó $f(x) = 2x + 1 \implies \lim_{x \to 1^+} f(x) = \lim_{x \to 1^+} (2x + 1) = 2(1) + 1 = 3$.
        \item Khi $x \to 1^-$ thì $x < 1$, do đó $f(x) = x^2 + 2 \implies \lim_{x \to 1^-} f(x) = \lim_{x \to 1^-} (x^2 + 2) = 1^2 + 2 = 3$.
        \item Vì $\lim_{x \to 1^+} f(x) = \lim_{x \to 1^-} f(x) = 3$ nên tồn tại giới hạn và $\lim_{x \to 1} f(x) = 3$.
    \end{itemize}
\end{itemize}

\noindent \textbf{d) Tổ chức thực hiện:}
\begin{itemize}[leftmargin=0.8cm]
    \item \textit{Chuyển giao nhiệm vụ:} Giáo viên trình bày các định lý; chiếu mô phỏng Desmos; yêu cầu học sinh làm Ví dụ 1 và Ví dụ 2 vào vở.
    \item \textit{Thực hiện nhiệm vụ:} Học sinh lắng nghe, ghi chép định lý, độc lập tính toán ví dụ trong 5 phút.
    \item \textit{Báo cáo, thảo luận:} Mời 2 học sinh lên bảng trình bày; học sinh lớp nhận xét; giáo viên chuẩn hóa các bước trình bày.
    \item \textit{Kết luận, nhận định:} Giáo viên nhấn mạnh: Đối với hàm số đa thức hoặc hàm phân thức xác định tại điểm $x_0$, ta chỉ cần thay trực tiếp $x = x_0$ vào công thức để tìm giới hạn.
\end{itemize}

\subsubsection*{3. Hoạt động 3: Luyện tập (13 phút)}

\noindent \textbf{a) Mục tiêu:}
Rèn luyện kỹ năng tính giới hạn một phía để xét sự tồn tại của giới hạn hàm số tại một điểm và tính giới hạn hữu hạn tại vô cực.

\noindent \textbf{b) Nội dung: Nhiệm vụ trên Phiếu học tập số 1}
\begin{pedabox}{Phiếu luyện tập số 1: Giới hạn một phía và giới hạn tại vô cực}
\begin{enumerate}[leftmargin=0.6cm]
    \item \textbf{Bài 1:} Cho hàm số $g(x) = \begin{cases} x + 3 & \text{khi } x > 2 \\ 2x - 1 & \text{khi } x \le 2 \end{cases}$. Tìm $\lim_{x \to 2^+} g(x)$, $\lim_{x \to 2^-} g(x)$ và xét xem có tồn tại $\lim_{x \to 2} g(x)$ không.
    \item \textbf{Bài 2:} Tính giới hạn tại vô cực: $A = \lim_{x \to +\infty} \dfrac{4x - 1}{2x + 3}$ và $B = \lim_{x \to -\infty} \dfrac{2x^2 + x - 1}{x^2 + 5}$.
\end{enumerate}
\end{pedabox}

\noindent \textbf{c) Sản phẩm: Lời giải chi tiết của học sinh}
\begin{enumerate}[leftmargin=0.6cm]
    \item \textbf{Bài 1:}
    \begin{itemize}
        \item Giới hạn phải: $\lim_{x \to 2^+} g(x) = \lim_{x \to 2^+} (x + 3) = 2 + 3 = 5$.
        \item Giới hạn trái: $\lim_{x \to 2^-} g(x) = \lim_{x \to 2^-} (2x - 1) = 2(2) - 1 = 3$.
        \item Nhận xét: Vì $\lim_{x \to 2^+} g(x) = 5 \neq \lim_{x \to 2^-} g(x) = 3$ nên \textbf{không tồn tại} giới hạn $\lim_{x \to 2} g(x)$.
    \end{itemize}
    \item \textbf{Bài 2:}
    \begin{itemize}
        \item $A = \lim_{x \to +\infty} \dfrac{4x - 1}{2x + 3} = \lim_{x \to +\infty} \dfrac{4 - \dfrac{1}{x}}{2 + \dfrac{3}{x}} = \dfrac{4 - 0}{2 + 0} = \dfrac{4}{2} = 2$.
        \item $B = \lim_{x \to -\infty} \dfrac{2x^2 + x - 1}{x^2 + 5} = \lim_{x \to -\infty} \dfrac{2 + \dfrac{1}{x} - \dfrac{1}{x^2}}{1 + \dfrac{5}{x^2}} = \dfrac{2 + 0 - 0}{1 + 0} = 2$.
    \end{itemize}
\end{enumerate}

\noindent \textbf{d) Tổ chức thực hiện:}
\begin{itemize}[leftmargin=0.8cm]
    \item \textit{Chuyển giao nhiệm vụ:} Học sinh làm việc cá nhân 6 phút, sau đó đổi vở chấm chéo theo cặp bàn.
    \item \textit{Thực hiện nhiệm vụ:} Học sinh tính toán, kiểm tra lại bằng nút CALC trên MTCT với $X = 2{,}000001$ và $X = 1{,}999999$.
    \item \textit{Báo cáo, thảo luận:} 2 học sinh đại diện lên bảng ghi lời giải; lớp thảo luận về kết luận "không tồn tại giới hạn".
    \item \textit{Kết luận, nhận định:} Giáo viên nhấn mạnh: Giới hạn trái và giới hạn phải chỉ cần lệch nhau dù một giá trị nhỏ nhất thì giới hạn hai phía đều không tồn tại.
\end{itemize}

\subsubsection*{4. Hoạt động 4: Vận dụng (7 phút)}

\noindent \textbf{a) Mục tiêu:}
Mô hình hóa bài toán kinh tế thực tiễn về hàm chi phí trung bình tiệm cận trong sản xuất công nghiệp.

\noindent \textbf{b) Nội dung: Bài toán chi phí trung bình dài hạn của doanh nghiệp}
\begin{pedabox}{Ứng dụng thực tiễn: Chi phí sản xuất trung bình khi sản lượng cực lớn}
Một công ty may mặc đầu tư dây chuyền sản xuất áo khoác mùa đông. Tổng chi phí sản xuất ra $x$ chiếc áo (tính bằng nghìn đồng) được cho bởi hàm số:
$$C(x) = 50\,000 + 120x \quad (x > 0),$$
trong đó $50\,000$ (nghìn đồng) là chi phí cố định (mặt bằng, máy móc) và $120$ (nghìn đồng) là chi phí sản xuất biến đổi cho mỗi chiếc áo.
\begin{enumerate}[leftmargin=0.6cm]
    \item Hãy viết công thức tính chi phí sản xuất trung bình cho mỗi chiếc áo: $\overline{C}(x) = \dfrac{C(x)}{x}$.
    \item Tính $\lim_{x \to +\infty} \overline{C}(x)$ và nêu ý nghĩa kinh tế của kết quả này đối với bài toán hiệu quả kinh tế theo quy mô (Economies of Scale).
\end{enumerate}
\end{pedabox}

\noindent \textbf{c) Sản phẩm: Bài giải của học sinh}
\begin{enumerate}[leftmargin=0.6cm]
    \item Chi phí trung bình: $\overline{C}(x) = \dfrac{50\,000 + 120x}{x} = \dfrac{50\,000}{x} + 120$ (nghìn đồng/chiếc áo).
    \item Ta có:
    $$\lim_{x \to +\infty} \overline{C}(x) = \lim_{x \to +\infty} \left(\dfrac{50\,000}{x} + 120\right) = 0 + 120 = 120 \text{ (nghìn đồng)}.$$
    \item \textbf{Ý nghĩa kinh tế:} Khi số lượng sản phẩm sản xuất ra vô cùng lớn ($x \to +\infty$), chi phí cố định ban đầu phân bổ cho mỗi chiếc áo tiến sát về 0. Chi phí sản xuất mỗi chiếc áo sẽ tiệm cận về chi phí biến đổi tối thiểu là $120\,000$ đồng, giúp hạ giá thành sản phẩm và nâng cao năng lực cạnh tranh.
\end{enumerate}

\noindent \textbf{d) Tổ chức thực hiện:}
Giáo viên đặt câu hỏi; học sinh thảo luận cặp đôi 3 phút và trả lời; giáo viên chốt kiến thức liên môn Toán học - Kinh tế học.

\newpage

\subsection*{TIẾT 2: GIỚI HẠN VÔ CỰC, DẠNG VÔ ĐỊNH VÀ BÀI TOÁN THỰC TIỄN (TIẾT 43 THEO PPCT)}

\subsubsection*{1. Hoạt động 1: Mở đầu (Khởi động) (5 phút)}

\noindent \textbf{a) Mục tiêu:}
Tạo tình huống dẫn dắt học sinh tiếp cận khái niệm giới hạn vô cực khi mẫu số tiến dần về 0.

\noindent \textbf{b) Nội dung: Khảo sát hàm số $f(x) = \dfrac{1}{x - 2}$ khi $x$ tiến sát 2}
\begin{pedabox}{Khởi động: Điều gì xảy ra khi chia cho một số cực kỳ nhỏ?}
Xét hàm số $f(x) = \dfrac{1}{x - 2}$ với $x > 2$.
\begin{itemize}[leftmargin=0.6cm]
    \item Cho $x$ tiến dần về $2$ từ bên phải: $x = 2{,}1 \implies f(x) = \dfrac{1}{0{,}1} = 10$.
    \item Khi $x = 2{,}01 \implies f(x) = \dfrac{1}{0{,}01} = 100$.
    \item Khi $x = 2{,}000001 \implies f(x) = 1\,000\,000$.
\end{itemize}
\textbf{Câu hỏi:} Khi $x$ càng tiến sát tới $2$ (với $x > 2$), hiệu số $(x - 2)$ là số dương tiến về $0$. Khi đó giá trị của phân số $f(x)$ trở nên như thế nào? Có bị chặn bởi một số cụ thể nào không?
\end{pedabox}

\noindent \textbf{c) Sản phẩm: Nhận xét của học sinh}
Học sinh nhận thấy giá trị $f(x)$ tăng lên vô hạn, lớn hơn bất kỳ số dương nào tùy ý. Ta nói giới hạn của $f(x)$ khi $x$ dần tới $2^+$ là dương vô cực ($+\infty$).

\noindent \textbf{d) Tổ chức thực hiện:}
Giáo viên trình chiếu các phép chia nhỏ dần; học sinh phát biểu cảm nhận; giáo viên chuẩn hóa khái niệm giới hạn vô cực.

\subsubsection*{2. Hoạt động 2: Hình thành kiến thức (18 phút)}

\noindent \textbf{a) Mục tiêu:}
\begin{itemize}[leftmargin=0.6cm]
    \item Nắm vững định nghĩa giới hạn vô cực của hàm số tại một điểm và các quy tắc tìm giới hạn của tích, thương liên quan đến vô cực.
    \item Làm chủ kỹ thuật khử dạng vô định $\dfrac{0}{0}$ bằng phân tích đa thức thành nhân tử và nhân lượng liên hợp.
    \item Tích hợp Năng lực số 5.2NC1b (thao tác CALC $X = x_0 + 10^{-9}$) và Năng lực AI 11.C2.1 (mô phỏng tải trọng tới hạn).
\end{itemize}

\noindent \textbf{b) Nội dung: Giới hạn vô cực và phương pháp khử dạng vô định}

\begin{pedabox}{1. Khái niệm giới hạn vô cực của hàm số tại một điểm}
\begin{itemize}[leftmargin=0.6cm]
    \item Ta nói hàm số $y = f(x)$ có giới hạn là $+\infty$ khi $x \to x_0$ nếu với mọi dãy số $(x_n)$ bất kì thỏa mãn $x_n \in K \setminus \{x_0\}$ và $x_n \to x_0$ thì ta đều có $f(x_n) \to +\infty$. Kí hiệu: $\lim_{x \to x_0} f(x) = +\infty$.
    \item Tương tự cho $\lim_{x \to x_0} f(x) = -\infty$ và các giới hạn một phía: $\lim_{x \to x_0^+} f(x) = \pm\infty$, $\lim_{x \to x_0^-} f(x) = \pm\infty$.
\end{itemize}
\end{pedabox}

\begin{scaffoldbox}{Hộp hỗ trợ sư phạm: Quy tắc xét dấu giới hạn của thương $\dfrac{f(x)}{g(x)}$}
Giả sử $\lim_{x \to x_0} f(x) = L \neq 0$ và $\lim_{x \to x_0} g(x) = 0$. Khi đó:
\begin{center}
\small
\begin{tabularx}{\textwidth}{|c|c|c|}
\hline
\textbf{Dấu của $L$} & \textbf{Dấu của $g(x)$ (trong lân cận $x_0$)} & \textbf{Giới hạn $\lim_{x \to x_0} \dfrac{f(x)}{g(x)}$} \\ \hline
$L > 0$ & $g(x) > 0$ & $+\infty$ \\ \hline
$L > 0$ & $g(x) < 0$ & $-\infty$ \\ \hline
$L < 0$ & $g(x) > 0$ & $-\infty$ \\ \hline
$L < 0$ & $g(x) < 0$ & $+\infty$ \\ \hline
\end{tabularx}
\end{center}
\end{scaffoldbox}

\begin{pedabox}{2. Phương pháp khử dạng vô định $\dfrac{0}{0}$ tại điểm $x \to x_0$}
Khi tính $\lim_{x \to x_0} \dfrac{P(x)}{Q(x)}$ mà $P(x_0) = 0$ và $Q(x_0) = 0$ (dạng vô định $\dfrac{0}{0}$):
\begin{itemize}[leftmargin=0.6cm]
    \item \textbf{Phương pháp 1 (Phân tích đa thức thành nhân tử):} Vì $P(x_0) = Q(x_0) = 0$ nên cả hai đa thức đều chia hết cho nhị thức $(x - x_0)$. Ta phân tích:
    $$\dfrac{P(x)}{Q(x)} = \dfrac{(x - x_0) \cdot P_1(x)}{(x - x_0) \cdot Q_1(x)} = \dfrac{P_1(x)}{Q_1(x)} \quad (\text{với } x \neq x_0).$$
    \item \textbf{Phương pháp 2 (Nhân lượng liên hợp khử căn thức):} Áp dụng các hằng đẳng thức:
    $$(\sqrt{A} - B)(\sqrt{A} + B) = A - B^2; \quad (\sqrt[3]{A} - B)(\sqrt[3]{A^2} + B\sqrt[3]{A} + B^2) = A - B^3.$$
\end{itemize}
\end{pedabox}

\begin{aibox}{Năng lực số 5.2NC1b: Kỹ thuật kiểm tra giới hạn trên MTCT Casio fx-580VN X}
Để tính $\lim_{x \to 2} \dfrac{x^2 - 5x + 6}{x - 2}$ và kiểm tra giới hạn một phía $\lim_{x \to 1^+} \dfrac{2x - 3}{x - 1}$:
\begin{enumerate}[leftmargin=0.6cm]
    \item \textbf{Khử dạng $\dfrac{0}{0}$:} Nhập \texttt{(X\textsuperscript{2} - 5X + 6) / (X - 2)} $\to$ Bấm \fbox{\textbf{CALC}} $\to$ Nhập \fbox{\textbf{2.000001}} $\to$ Kết quả hiển thị: \texttt{-0.999999}. Dự đoán kết quả bằng $-1$.
    \item \textbf{Giới hạn vô cực:} Nhập \texttt{(2X - 3) / (X - 1)} $\to$ Bấm \fbox{\textbf{CALC}} $\to$ Nhập \fbox{\textbf{1 + 10\textsuperscript{-9}}} $\to$ Kết quả hiển thị: \texttt{-1 $\times$ 10\textsuperscript{9}} (số âm rất lớn). Kết luận giới hạn là $-\infty$.
\end{enumerate}
\end{aibox}

\noindent \textbf{c) Sản phẩm: Lời giải các ví dụ mẫu}
\textbf{Ví dụ 1 (Khử dạng $\dfrac{0}{0}$ đa thức):} Tính $L_1 = \lim_{x \to 2} \dfrac{x^2 - 5x + 6}{x - 2}$.
\begin{itemize}[leftmargin=0.6cm]
    \item Khi thay $x = 2$: Tử $= 2^2 - 10 + 6 = 0$, Mẫu $= 2 - 2 = 0$ (dạng $\dfrac{0}{0}$).
    \item Phân tích tử số: $x^2 - 5x + 6 = (x - 2)(x - 3)$.
    \item Ta có:
    $$L_1 = \lim_{x \to 2} \dfrac{(x - 2)(x - 3)}{x - 2} = \lim_{x \to 2} (x - 3) = 2 - 3 = -1.$$
\end{itemize}

\textbf{Ví dụ 2 (Khử dạng $\dfrac{0}{0}$ chứa căn thức):} Tính $L_2 = \lim_{x \to 3} \dfrac{\sqrt{x + 1} - 2}{x - 3}$.
\begin{itemize}[leftmargin=0.6cm]
    \item Khi thay $x = 3$: Tử $= \sqrt{4} - 2 = 0$, Mẫu $= 0$ (dạng $\dfrac{0}{0}$).
    \item Nhân cả tử và mẫu với lượng liên hợp $(\sqrt{x + 1} + 2)$:
    $$L_2 = \lim_{x \to 3} \dfrac{(\sqrt{x + 1} - 2)(\sqrt{x + 1} + 2)}{(x - 3)(\sqrt{x + 1} + 2)} = \lim_{x \to 3} \dfrac{(x + 1) - 4}{(x - 3)(\sqrt{x + 1} + 2)} = \lim_{x \to 3} \dfrac{x - 3}{(x - 3)(\sqrt{x + 1} + 2)}.$$
    \item Rút gọn $(x - 3)$:
    $$L_2 = \lim_{x \to 3} \dfrac{1}{\sqrt{x + 1} + 2} = \dfrac{1}{\sqrt{3 + 1} + 2} = \dfrac{1}{2 + 2} = \dfrac{1}{4}.$$
\end{itemize}

\noindent \textbf{d) Tổ chức thực hiện:}
\begin{itemize}[leftmargin=0.8cm]
    \item \textit{Chuyển giao nhiệm vụ:} Giáo viên hướng dẫn nhận dạng $\dfrac{0}{0}$; thị phạm cách nhân lượng liên hợp; học sinh thực hành theo dõi.
    \item \textit{Thực hiện nhiệm vụ:} Học sinh ghi lại quy tắc nhân liên hợp; tự làm Ví dụ 1 và 2 vào vở.
    \item \textit{Báo cáo, thảo luận:} Giáo viên kiểm tra vở của một số học sinh yếu; gọi 1 bạn đọc các bước biến đổi căn thức.
    \item \textit{Kết luận, nhận định:} Giáo viên nhấn mạnh: Mục tiêu tối hậu của việc phân tích hay nhân liên hợp ở dạng $\dfrac{0}{0}$ là làm xuất hiện và \textbf{triệt tiêu nhân tử chung $(x - x_0)$}.
\end{itemize}

\subsubsection*{3. Hoạt động 3: Luyện tập (15 phút)}

\noindent \textbf{a) Mục tiêu:}
Củng cố kỹ năng tính giới hạn một phía vô cực và khử thành thạo dạng vô định $\dfrac{0}{0}$ chứa căn thức bậc hai.

\noindent \textbf{b) Nội dung: Nhiệm vụ trên Phiếu học tập số 2}
\begin{pedabox}{Phiếu luyện tập số 2: Giới hạn vô cực và khử dạng vô định}
Tính các giới hạn sau bằng tự luận và đối chiếu kết quả MTCT:
\begin{enumerate}[leftmargin=0.6cm]
    \item \textbf{Bài 1 (Giới hạn một phía vô cực):} Tính $J_1 = \lim_{x \to 1^+} \dfrac{2x - 3}{x - 1}$ và $J_2 = \lim_{x \to 1^-} \dfrac{2x - 3}{x - 1}$.
    \item \textbf{Bài 2 (Khử dạng $\dfrac{0}{0}$ đa thức):} Tính $J_3 = \lim_{x \to -1} \dfrac{x^3 + 1}{x^2 - 1}$.
    \item \textbf{Bài 3 (Khử dạng $\dfrac{0}{0}$ căn thức):} Tính $J_4 = \lim_{x \to 0} \dfrac{\sqrt{1 + 3x} - 1}{x}$.
\end{enumerate}
\end{pedabox}

\noindent \textbf{c) Sản phẩm: Lời giải chi tiết}
\begin{enumerate}[leftmargin=0.6cm]
    \item \textbf{Bài 1:}
    \begin{itemize}
        \item Xét $J_1$: Ta có $\lim_{x \to 1^+} (2x - 3) = 2(1) - 3 = -1 < 0$. Khi $x \to 1^+$ thì $x > 1 \implies x - 1 > 0$. Do đó $J_1 = -\infty$.
        \item Xét $J_2$: Ta có $\lim_{x \to 1^-} (2x - 3) = -1 < 0$. Khi $x \to 1^-$ thì $x < 1 \implies x - 1 < 0$. Do đó $J_2 = +\infty$.
    \end{itemize}
    \item \textbf{Bài 2:} Khi $x \to -1$: Cả tử và mẫu bằng 0. Áp dụng hằng đẳng thức:
    $$J_3 = \lim_{x \to -1} \dfrac{(x + 1)(x^2 - x + 1)}{(x + 1)(x - 1)} = \lim_{x \to -1} \dfrac{x^2 - x + 1}{x - 1} = \dfrac{(-1)^2 - (-1) + 1}{-1 - 1} = \dfrac{3}{-2} = -\dfrac{3}{2}.$$
    \item \textbf{Bài 3:} Nhân liên hợp với $(\sqrt{1 + 3x} + 1)$:
    $$J_4 = \lim_{x \to 0} \dfrac{(\sqrt{1 + 3x} - 1)(\sqrt{1 + 3x} + 1)}{x(\sqrt{1 + 3x} + 1)} = \lim_{x \to 0} \dfrac{(1 + 3x) - 1}{x(\sqrt{1 + 3x} + 1)} = \lim_{x \to 0} \dfrac{3x}{x(\sqrt{1 + 3x} + 1)}.$$
    Rút gọn $x$:
    $$J_4 = \lim_{x \to 0} \dfrac{3}{\sqrt{1 + 3x} + 1} = \dfrac{3}{\sqrt{1} + 1} = \dfrac{3}{2}.$$
\end{enumerate}

\noindent \textbf{d) Tổ chức thực hiện:}
\begin{itemize}[leftmargin=0.8cm]
    \item \textit{Chuyển giao nhiệm vụ:} 3 nhóm học sinh thực hiện 3 bài tập trên giấy nháp; chuẩn bị báo cáo lời giải.
    \item \textit{Thực hiện nhiệm vụ:} Học sinh giải bài độc lập; giáo viên nhắc nhở học sinh chú ý hằng đẳng thức $a^3 + b^3 = (a+b)(a^2 - ab + b^2)$.
    \item \textit{Báo cáo, thảo luận:} Gọi 3 học sinh lên bảng hoàn thiện lời giải; lớp đối chiếu với màn hình MTCT.
    \item \textit{Kết luận, nhận định:} Giáo viên nhận xét, cho điểm học sinh có lời giải rõ ràng, mạch lạc.
\end{itemize}

\subsubsection*{4. Hoạt động 4: Vận dụng \& Năng lực AI (7 phút)}

\noindent \textbf{a) Mục tiêu:}
Tích hợp Năng lực AI 11.C2.1 và rèn luyện tư duy phản biện lỗi sai phổ biến của AI khi tính toán giới hạn hàm số phân nhánh hoặc không tồn tại.

\noindent \textbf{b) Nội dung: Phản biện câu trả lời của AI về giới hạn hàm chứa giá trị tuyệt đối}
\begin{aibox}{Năng lực AI: Phản biện sai sót của Chatbot AI khi giải toán giải tích}
\textbf{Tình huống thực tế:} Một bạn học sinh đặt câu hỏi cho ứng dụng AI trên điện thoại:
\begin{quote}
\textit{«Hãy tính giới hạn: $L = \lim_{x \to 1} \dfrac{x - 1}{|x - 1|}$.»}
\end{quote}
\textbf{Câu trả lời của AI:}
\begin{quote}
\textit{«Vì biểu thức có dạng $\dfrac{0}{0}$ nên ta rút gọn tử và mẫu: $\dfrac{x - 1}{|x - 1|} = 1$. Do đó giới hạn $L = 1$!»}
\end{quote}
\textbf{Nhiệm vụ của học sinh:}
\begin{enumerate}[leftmargin=0.6cm]
    \item Chỉ ra sai lầm nghiêm trọng trong suy luận của AI.
    \item Trình bày lại lời giải đúng và đưa ra kết luận chính xác.
    \item Nêu ý nghĩa ứng dụng AI (Mã 11.C2.1): Tại sao các mô hình AI tính toán kỹ thuật cần thuật toán kiểm tra giới hạn vô cực để cảnh báo tải trọng tới hạn của cầu đường?
\end{enumerate}
\end{aibox}

\noindent \textbf{c) Sản phẩm: Lời phản biện và phân tích của học sinh}
\begin{enumerate}[leftmargin=0.6cm]
    \item \textbf{Lỗi sai của AI:} AI đã tự ý bỏ dấu giá trị tuyệt đối mà không xét dấu của biểu thức $(x - 1)$, ngộ nhận rằng $|x - 1| = x - 1$ với mọi $x$.
    \item \textbf{Lời giải chuẩn xác:}
    \begin{itemize}
        \item Giới hạn phải: Khi $x \to 1^+$ thì $x > 1 \implies x - 1 > 0 \implies |x - 1| = x - 1$.
        Do đó: $\lim_{x \to 1^+} \dfrac{x - 1}{|x - 1|} = \lim_{x \to 1^+} \dfrac{x - 1}{x - 1} = 1$.
        \item Giới hạn trái: Khi $x \to 1^-$ thì $x < 1 \implies x - 1 < 0 \implies |x - 1| = -(x - 1)$.
        Do đó: $\lim_{x \to 1^-} \dfrac{x - 1}{|x - 1|} = \lim_{x \to 1^-} \dfrac{x - 1}{-(x - 1)} = -1$.
        \item Vì $\lim_{x \to 1^+} \dfrac{x - 1}{|x - 1|} = 1 \neq \lim_{x \to 1^-} \dfrac{x - 1}{|x - 1|} = -1$ nên \textbf{KHÔNG TỒN TẠI} giới hạn $\lim_{x \to 1} \dfrac{x - 1}{|x - 1|}$.
    \end{itemize}
    \item \textbf{Ứng dụng AI kỹ thuật (11.C2.1):} Trong xây dựng cầu treo hoặc nhà cao tầng, khi gió bão hoặc đoàn xe tạo ra tần số dao động tiệm cận tần số riêng của công trình ($f \to f_0$), biên độ rung lắc $A(f) \to +\infty$ (hiện tượng cộng hưởng). Các thuật toán AI mô phỏng phải cảnh báo sớm giới hạn vô cực này để kỹ sư kịp thời gia cố kết cấu.
\end{enumerate}

\noindent \textbf{d) Tổ chức thực hiện:}
Giáo viên trình chiếu đoạn văn bản giải sai của AI; học sinh phát hiện lỗi sai và giơ tay giải thích; giáo viên chốt lại bài học và dặn dò bài tập về nhà.

\section*{IV. PHỤ LỤC HỌC LIỆU BỔ TRỢ}

\subsection*{1. Phiếu học tập phân tầng bậc thang (Worksheet)}

\noindent \textbf{A. PHẦN TRẮC NGHIỆM KHÁCH QUAN (Nhận biết - Thông hiểu)}
\begin{enumerate}[leftmargin=0.6cm]
    \item Giá trị của giới hạn $\lim_{x \to 2} (x^2 - 3x + 5)$ bằng:
    \begin{itemize}[leftmargin=0.6cm]
        \item[A.] $1$. \quad \textbf{B.} $3$. \quad C. $5$. \quad D. $7$.
    \end{itemize}
    \item Cho $\lim_{x \to 1} f(x) = 2$ và $\lim_{x \to 1} g(x) = -3$. Khi đó $\lim_{x \to 1} [2f(x) - g(x)]$ bằng:
    \begin{itemize}[leftmargin=0.6cm]
        \item[A.] $1$. \quad B. $-7$. \quad \textbf{C.} $7$. \quad D. $4$.
    \end{itemize}
    \item Giá trị của $\lim_{x \to 2^+} \dfrac{1}{x - 2}$ bằng:
    \begin{itemize}[leftmargin=0.6cm]
        \item[\textbf{A.}] $+\infty$. \quad B. $-\infty$. \quad C. $0$. \quad D. Không tồn tại.
    \end{itemize}
    \item Giới hạn $\lim_{x \to 1} \dfrac{x^2 - 1}{x - 1}$ bằng:
    \begin{itemize}[leftmargin=0.6cm]
        \item[A.] $0$. \quad B. $1$. \quad \textbf{C.} $2$. \quad D. $+\infty$.
    \end{itemize}
\end{enumerate}

\noindent \textbf{B. PHẦN TỰ LUẬN PHÂN TẦNG BẬC THANG}
\begin{itemize}[leftmargin=0.6cm]
    \item \textbf{Tầng 1 (Cơ bản dành cho HS trung bình - yếu):}
    Tính các giới hạn sau bằng phương pháp thay số trực tiếp hoặc phân tích nhân tử đơn giản:
    \begin{enumerate}[leftmargin=0.5cm]
        \item $\lim_{x \to 3} \dfrac{2x + 1}{x - 1}$.
        \item $\lim_{x \to 1} \dfrac{x^2 - 3x + 2}{x - 1}$.
    \end{enumerate}

    \item \textbf{Tầng 2 (Thông hiểu):}
    \begin{enumerate}[leftmargin=0.5cm]
        \item Tính giới hạn khử dạng vô định chứa căn thức: $\lim_{x \to 4} \dfrac{\sqrt{2x + 1} - 3}{x - 4}$.
        \item Cho hàm số $h(x) = \begin{cases} x^2 + ax + 1 & \text{khi } x \ge 2 \\ 2x + 3 & \text{khi } x < 2 \end{cases}$. Tìm tham số $a$ để tồn tại giới hạn $\lim_{x \to 2} h(x)$.
    \end{enumerate}

    \item \textbf{Tầng 3 (Vận dụng cao):}
    \begin{enumerate}[leftmargin=0.5cm]
        \item Tính giới hạn chứa hai căn thức bậc hai và bậc ba: $\lim_{x \to 0} \dfrac{\sqrt{1 + 2x} - \sqrt[3]{1 + 3x}}{x^2}$.
        \item Một thấu kính hội tụ có tiêu cự $f = 20\text{ cm}$. Khoảng cách từ vật đến thấu kính là $d$ ($d > 20$). Khoảng cách từ ảnh đến thấu kính được xác định bởi công thức $d' = \dfrac{20d}{d - 20}$. Khảo sát giới hạn $\lim_{d \to 20^+} d'$ và $\lim_{d \to +\infty} d'$, giải thích hiện tượng quang học thu được.
    \end{enumerate}
\end{itemize}

\subsection*{2. Bảng Rubric đánh giá năng lực học sinh trong bài học}

\begin{center}
\small
\begin{tabularx}{\textwidth}{|p{2.8cm}|X|X|X|X|}
\hline
\textbf{Tiêu chí} & \textbf{Chưa đạt (Mức 1)} & \textbf{Đạt (Mức 2)} & \textbf{Khá (Mức 3)} & \textbf{Tốt (Mức 4)} \\ \hline
\textbf{Nhận thức khái niệm \& Giới hạn một phía} & Chưa phân biệt được $x \to x_0$ và $x = x_0$; lúng túng khi xét giới hạn trái/phải. & Nêu được định nghĩa giới hạn một phía; tính được hàm đa thức cơ bản. & Hiểu sâu điều kiện tồn tại giới hạn hai phía; tính đúng giới hạn hàm phân nhánh. & Nắm vững bản chất giải tích; phân tích thành thạo hàm chứa giá trị tuyệt đối phức tạp. \\ \hline
\textbf{Kỹ năng khử dạng vô định $\dfrac{0}{0}$} & Không nhận ra dạng vô định $\dfrac{0}{0}$; phân tích nhân tử sai dấu. & Phân tích được tam thức bậc hai; rút gọn được nhân tử chung $(x - x_0)$. & Khử thành thạo dạng vô định chứa căn bậc hai bằng lượng liên hợp. & Xử lý điêu luyện bài toán thêm bớt số hạng khử đồng thời căn bậc hai và căn bậc ba. \\ \hline
\textbf{Giới hạn vô cực \& Mô hình thực tế} & Nhầm lẫn giữa giới hạn bằng 0 và giới hạn vô cực; xét sai dấu quy tắc thương. & Nhớ bảng quy tắc dấu thương; tính được giới hạn một phía ra vô cực đơn giản. & Vận dụng tốt vào bài toán kinh tế chi phí trung bình và bài toán quang học thấu kính. & Thiết lập mô hình giải tích xuất sắc; giải thích cặn kẽ hiện tượng tiệm cận trong đời sống. \\ \hline
\textbf{Năng lực số \& Phản biện AI} & Chưa biết dùng CALC kiểm tra nghiệm; tin tưởng tuyệt đối vào lời giải sai của AI. & Sử dụng được MTCT tính giá trị xấp xỉ; phát hiện được AI làm sai nhưng chưa sửa được. & Thao tác thuần thục CALC $X = x_0 \pm 10^{-9}$; phản biện sắc bén lỗi bỏ trị tuyệt đối của AI. & Làm chủ công cụ số; kết nối sâu sắc giới hạn vô cực với mô phỏng tải trọng tới hạn của AI kỹ thuật. \\ \hline
\end{tabularx}
\end{center}

\end{document}
